\SourceLeaf{444}%
folgt der wichtige Lehrsatz, dass
\begin{quote}
»bei allen Flächen, die sich in eine Ebene abwickeln lassen, das Krümmungsmaß überall verschwindet«
\end{quote}
\noindent oder dass
\[
\left(\frac{\partial\partial z}{\partial x.\partial y}\right)^{2}-\left(\frac{\partial\partial z}{\partial x^{2}}\right)\left(\frac{\partial\partial z}{\partial y^{2}}\right)=0
\]
gilt. Dieses Kriterium wird sonst aus anderen Gründen hergeleitet, wenn auch, wie uns scheint, mit geringerer Deutlichkeit, als man wünschen würde. Es ist klar, dass bei allen solchen Flächen die Zenite sämtlicher Punkte kein Gebiet ausfüllen können und daher alle auf einer Linie liegen.

\Needspace{5\baselineskip}
\begin{center}17.\end{center}
\noindent\hspace*{1.2em}Von einem gegebenen Punkt auf einer krummen Fläche lassen wir unzählige kürzeste Linien auf dieser Fläche ausgehen. Sie sollen sich durch den Winkel unterscheiden, den ihr erstes Element mit dem ersten Element einer \textls[100]{bestimmten} kürzesten Linie bildet. Dieser Winkel soll $\theta$ heißen. Ferner sei $s$ die [vom gegebenen Punkt aus gemessene] Länge eines Stücks einer solchen kürzesten Linie, und deren Endpunkte sollen die Koordinaten $x,y,.z$ haben. Da $\theta$ und $s$ also zu einem ganz bestimmten Punkt der krummen Fläche gehören, kann man $x,y,z$ als Funktionen von $\theta$ und $s$ betrachten. Der Richtung des Elements von $s$ entspreche auf der Kugel der Punkt $\lambda$ mit den Koordinaten $\xi,\eta,\zeta$, sodass
\[
\xi=\frac{\partial x}{\partial s},\qquad\eta=\frac{\partial y}{\partial s},\qquad\zeta=\frac{\partial z}{\partial s}
\]
gilt.

Die Endpunkte aller kürzesten Linien derselben Länge $s$ entsprechen einer krummen Linie, deren Länge $t$ heißen mag. Offenbar kann man $t$ als Funktion von $s$ und $\theta$ betrachten. Entspricht der Richtung des Elements von $t$ auf der Kugel der Punkt $\lambda'$ mit den Koordinaten $\xi',\eta',\zeta'$, so gilt
\[
\xi'.\frac{\partial t}{\partial\theta}=\frac{\partial x}{\partial\theta},\qquad\eta'.\frac{\partial t}{\partial\theta}=\frac{\partial y}{\partial\theta},\qquad\zeta'.\frac{\partial t}{\partial\theta}=\frac{\partial z}{\partial\theta}
\]
Somit gilt
